% TOP_PROBLEM: {"id": 30006780, "title": "Polyakov Conjecture: Absence of a Phase Transition for Planar O(n) Models with n >= 3", "category_id": 16, "category": {"id": 16, "name": "physics", "display_name": "Mathematical Physics", "description": "Problems at the intersection of mathematics and physics.", "slug": "physics", "order_index": 16, "created_at": "2026-07-31T15:26:25.671Z"}, "status": "open"}
% FIRST_POSTED: 2026-09-27
% !TeX program = lualatex
% Compile twice: lualatex open_problems_500.tex
% All references and prose are embedded; no BibTeX database or external inputs.
\documentclass[11pt,a4paper]{article}
\usepackage[margin=24mm,headheight=14pt]{geometry}
\usepackage{fontspec}
\setmainfont{Latin Modern Roman}
\setsansfont{Latin Modern Sans}
\setmonofont{Latin Modern Mono}
\usepackage{amsmath,amssymb,mathtools}
\usepackage{unicode-math}
\setmathfont{Latin Modern Math}
\usepackage{microtype}
\usepackage{enumitem,xurl,needspace}
\usepackage[dvipsnames]{xcolor}
\usepackage{hyperref}
\hypersetup{unicode=true,colorlinks=true,linkcolor=MidnightBlue,urlcolor=MidnightBlue,
 pdftitle={500 Mathematical Problem Statements},pdfauthor={Source-annotated compilation},bookmarksnumbered=true}
\usepackage{bookmark}
\usepackage{tocloft}
\setlength{\cftsecnumwidth}{3em}
\cftsetpnumwidth{2.5em}
\cftsetrmarg{4em}
\usepackage{titlesec}
\titleformat{\section}{\Large\bfseries\raggedright}{\thesection}{0.7em}{}
\usepackage{fancyhdr}
\pagestyle{fancy}\fancyhf{}
\fancyhead[L]{\small\sffamily Mathematical problem statements}
\fancyhead[R]{\small Source edition 22}
\fancyfoot[C]{\thepage}
\setlength{\parindent}{0pt}
\setlength{\parskip}{5pt plus 1pt}
\setlength{\emergencystretch}{3em}
\setcounter{tocdepth}{1}
\setcounter{secnumdepth}{1}
\setlist[itemize]{leftmargin=3em,itemsep=5pt,topsep=4pt,parsep=0pt}
\allowdisplaybreaks
\newcommand{\N}{\mathbb{N}}
\newcommand{\Z}{\mathbb{Z}}
\newcommand{\Q}{\mathbb{Q}}
\newcommand{\R}{\mathbb{R}}
\newcommand{\C}{\mathbb{C}}
\newcommand{\F}{\mathbb{F}}
\newcommand{\A}{\mathbb{A}}
\newcommand{\PP}{\mathbb P}
\newcommand{\E}{\mathbb{E}}
\newcommand{\eps}{\varepsilon}
\newcommand{\dd}{\,\mathrm d}
\newcommand{\sref}[2]{\hyperlink{source:#1:#2}{[S#2]}}
\newcommand{\eref}[2]{\hyperlink{extra:#1:#2}{[E#2]}}
% Turn the next switch off for a shorter reading copy. The quotations preserve
% every exactTarget field, including qualifications omitted from a short summary.
\newif\ifshowcatalogue
\showcataloguetrue
\newcommand{\cataloguescope}[1]{\ifshowcatalogue
 \paragraph{Exact scope in the supplied catalogue.}
 \begin{quote}\small #1\end{quote}\fi}

% Compile twice with: lualatex 30006780.tex
\numberwithin{equation}{section}
\hypersetup{pdftitle={Research attempt -- problem 30006780},pdfauthor={Research notebook; original compilation retained}}
\fancyhead[L]{\small\sffamily Research notebook}
\fancyhead[R]{\small\sffamily Problem 30006780 | 27 September 2026}
\begin{document}
\setcounter{section}{0}
% ================================================================
% problemId: 30006780
% Canonical problem ID: 30006780
% exactTarget SHA-256: 1cf447cbe644bd586a467045b264bcc37f6d18eb06cbda2969a2553578ca629e
\clearpage
\section*{\texorpdfstring{Polyakov Conjecture: Absence of a Phase Transition for Planar O(n) Models with n >= 3}{Polyakov Conjecture: Absence of a Phase Transition for Planar O(n) Models with n >= 3}}\label{30006780}
\subsection{Definitions and mathematical statement}
Place spins $\sigma_x\in S^{n-1}\subset\R^n$ at $x\in\Z^2$. The nearest-neighbor ferromagnetic Gibbs weight in a finite region is proportional to
$\exp(\beta\sum_{\langle x,y\rangle}\sigma_x\cdot\sigma_y)$ relative to product uniform spherical measure. Two-point correlation is $\langle\sigma_0\cdot\sigma_x\rangle_\beta$ in the infinite-volume Gibbs convention of the source. For each $n\ge3$ and finite $\beta>0$, the target is exponential decay:
\[
 |\langle\sigma_0\cdot\sigma_x\rangle_\beta|\le C(n,\beta)e^{-c(n,\beta)|x|},
 \qquad c(n,\beta)>0.
\]
The constants need not be uniform as temperature tends to zero. Absence of spontaneous magnetization alone is weaker than the claimed exponential bound.
\par\smallskip\textit{Formulation and concept references:} \sref{30006780}{1}.
\cataloguescope{For every integer n >= 3 and every finite inverse temperature beta > 0, prove that two-point correlations in the nearest-neighbor O(n) spin model on Z\textasciicircum{}2 decay exponentially with distance.}
\subsection{Short English statement}
Do planar O(n) spin correlations decay exponentially at every positive temperature once the spin dimension is at least three?
% BEGIN RESEARCH ADDITION ATTEMPT-187
\subsection{Research attempt}

\paragraph{Current frontier.}
Absence of spontaneous magnetization in two dimensions is weaker than a mass gap. Aru--Garban--Sep\'ulveda formulate the $n\ge3$ exponential-decay problem and study a percolative mechanism related to vector-valued Gaussian free fields; their results must not be read as a proof of exponential decay for every finite inverse temperature \sref{30006780}{1}, Section~1.1. We first reconstruct a high-temperature contraction, then isolate exactly what a scale-dependent improvement would need. The distinction from the planar $O(2)$ model is essential: that model has a rigorously established low-temperature nonexponential regime \eref{30006780}{1}.

\paragraph{Attack route.}
Single-site mixing can give an explicit high-temperature criterion. A coarse-grained version could in principle extend it to all finite $\beta$, but only if an $n$-dependent mechanism makes a block-to-block influence kernel strictly subcritical. Alternatively one could prove a suitable finite-volume susceptibility criterion directly. We investigate the contraction route; neither a percolation analogy nor absence of magnetization is substituted for the missing contraction estimate.

\paragraph{Attempt.}
At a site with neighbor sum $h$, the conditional spin distribution is
\[
 d\nu_h(s)=Z(h)^{-1}e^{\beta s\cdot h}\,d\omega(s),\qquad s\in S^{n-1}.
\]
When one neighboring spin changes from $u$ to $v$, the oscillation of the unnormalized log likelihood ratio is at most
\[
 \sup_s\beta s\cdot(u-v)-\inf_s\beta s\cdot(u-v)
 =2\beta|u-v|\le4\beta.
\]
If two probability densities have log likelihood ratio of oscillation at most $L$, their total variation distance is at most $\tanh(L/4)$. Here is a proof of the constant: for $r=dP/dQ$ with $a\le r\le b$ and $b/a\le e^L$, $\E_Qr=1$ gives
\[
 \tfrac12\E_Q|r-1|\le\frac{(1-a)(b-1)}{b-a}.
\]
Maximizing over $a\le1\le b$, $b/a\le e^L$, gives the symmetric extremum $a=e^{-L/2}$, $b=e^{L/2}$ and the asserted bound. Thus each of the four neighbors has influence at most $t=\tanh\beta$.

For completeness, the resulting comparison has a probabilistic interpretation. Couple two finite-volume heat-bath systems by maximal couplings. A disagreement at $x$ is propagated to a neighbor with coefficient at most $t$, so iteration gives a sum over nearest-neighbor walks. If
\begin{equation}\label{eq187:smallbeta}
 q=4\tanh\beta<1,
\end{equation}
that sum converges and the influence between sites at graph distance $r$ is at most
\[
 \sum_{j\ge r}4^jt^j=\frac{q^r}{1-q}.
\]
The Dobrushin comparison theorem, whose proof uses these coupled finite-volume inequalities, makes this argument rigorous \eref{30006780}{2}. One solves the finite-volume disagreement inequalities $d\le b+Cd$, with $C_{xy}=t$ on neighboring sites, and using $(I-C)^{-1}=\sum_{j\ge0}C^j$. Exhausting the lattice makes boundary influences tend to zero and yields uniqueness. Rotation invariance then gives $\langle\sigma_x\rangle=0$.

Conditioning on the spin at the origin and applying the comparison to each bounded coordinate function gives the conservative bound
\begin{equation}\label{eq187:highT}
 |\langle\sigma_0\cdot\sigma_x\rangle|
 \le \frac{2n}{1-q}q^{|x|_1}\qquad(x\ne0).
\end{equation}
Indeed, the conditional expectation of each coordinate at $x$ changes by at most twice the influence bound as the spin at zero varies; its unconditional mean is zero. Multiply by the corresponding coordinate at zero, bounded by one, and sum over the $n$ coordinates. The factor $2n$ is not optimized. This recovers a high-temperature special case for $\beta<\operatorname{arctanh}(1/4)$, not the full assertion.

\textbf{Finite-scale lemma.} Let $g:\Z^2\to[0,1]$, and suppose for some integer $L\ge1$ and nonnegative kernel $K$ supported on $0<|z|_1\le L$ that
\begin{equation}\label{eq187:block}
 g(x)\le\sum_zK(z)g(x-z)\quad(|x|_1>L),\qquad
 \sum_zK(z)=q_L<1.
\end{equation}
Then, with $m=\max\{0,\lfloor |x|_1/L\rfloor-1\}$,
\[
 g(x)\le q_L^m.
\]
To prove it, substitute the inequality $m$ times. Before the last substitution every reachable site still has norm greater than $L$, by the triangle inequality. The sum of all path weights is exactly $q_L^m$, while the terminal value is at most one. If $q_L=0$ the function vanishes outside the ball; otherwise this is exponential decay with rate $-L^{-1}\log q_L$, up to a harmless prefactor.

Applying this lemma to the absolute correlation would prove the stated target if, for every $n\ge3$ and finite $\beta$, one could construct an actual correlation comparison kernel satisfying \eqref{eq187:block}. The lemma accommodates a correlation length growing arbitrarily rapidly with $\beta$; it does not demand an unphysical temperature-uniform mass gap.

\paragraph{Stress test.}
The single-site estimate is independent of the distinction $n=2$ versus $n\ge3$. Therefore iterating it without new geometric information cannot explain the conjectured distinction. Worst-case block boundary influences may become larger, not smaller: replacing a block by a ``renormalized spin'' and assuming a smaller effective coupling is an unproved step. The finite-scale lemma is only a sufficient analytic criterion; no Simon--Lieb-type inequality with subcritical total weight at all temperatures has been proved here. Mermin--Wagner's conclusion does not provide that kernel, and slow decay to zero would not suffice. The low-temperature $O(2)$ results in \eref{30006780}{1} actively falsify any dimension-insensitive argument purporting to establish the full all-temperature conclusion.

\paragraph{Outcome.}
\textbf{Outcome: Exploratory approach with unresolved gap.}
We proved an explicit high-temperature estimate and a finite-scale contraction lemma, both with transparent constants and quantifiers. They are elementary reconstructions, not a new mass-gap theorem. The missing step is an $n\ge3$-specific low-temperature estimate that produces a subcritical comparison kernel at some finite scale for every fixed $\beta$. No such estimate is established.

% END RESEARCH ADDITION ATTEMPT-187
\subsection{Sources}
\begin{itemize}\raggedright
\item[\textnormal{[S1]}]\hypertarget{source:30006780:1}{} J. Aru, C. Garban, and A. Sepulveda, Percolation for 2D Classical Heisenberg Model and Exit Sets of Vector Valued GFF, Section 1.1, CMP 406 (2025).
\url{https://link.springer.com/article/10.1007/s00220-024-05208-y}.
{\small\textit{Catalogue formulation source.}}
% BEGIN RESEARCH ADDITION REFERENCES-187
\item[\textnormal{[E1]}]\hypertarget{extra:30006780:1}{} J\"urg Fr\"ohlich and Thomas Spencer, \emph{The Kosterlitz--Thouless transition in two-dimensional abelian spin systems and the Coulomb gas}, Communications in Mathematical Physics 81 (1981), 527--602; low-temperature planar rotator correlations.
\url{https://doi.org/10.1007/BF01208273}.
{\small\textit{Research reference; consulted 27 September 2026.}}
\item[\textnormal{[E2]}]\hypertarget{extra:30006780:2}{} R. L. Dobrushin, \emph{The Description of a Random Field by Means of Conditional Probabilities and Conditions of Its Regularity}, Theory of Probability and Its Applications 13 (1968), no. 2, 197--224. Single-site influence contraction and comparison.
\url{https://doi.org/10.1137/1113026}.
{\small\textit{Research reference; consulted 27 September 2026.}}
% END RESEARCH ADDITION REFERENCES-187
\end{itemize}

\end{document}
